\documentclass{report}
\usepackage{amsfonts, amsmath}

\newcommand\N{{\mathbb N}}
\newcommand\Q{{\mathbb Q}}
\newcommand\R{{\mathbb R}}
\newcommand\Z{{\mathbb Z}}

\pagestyle{empty}
\begin{document}
\large
\centerline{\Huge \bf Analisi Matematica II}
\vskip.2cm
\centerline{\Huge  Terzo Appello (11-09-2002)}
\renewcommand{\thefootnote}{ } 
\footnote{\sl Avete 3 ore di tempo. Ogni esercizio
vale sei punti. La sufficienza si ottiene con un punteggio $\geq$
18. {\bf Solo le risposte chiaramente giustificate saranno prese in
considerazione.} {\scshape Elaborati illeggibili o confusi verranno ignorati.}
{\sffamily La
sintesi sar\`a particolarmente apprezzata.}}
\renewcommand{\thefootnote}{\arabic{footnote}} 

\vskip1cm
\begin{enumerate}
\item Un banca vi offre il seguente piano di investimento: 
se depositate un capitale di
$x$ euro, vincolato per sei anni, allo scadere del sesto anno la banca vi
restituisce $1.4 x$ euro.  Tenendo conto che la banca paga gli interessi
su base bimestrale, quale \`e l'interesse annuo?
\item Sia $f$ derivabile due volte in ogni punto. Si dimostri che
\[
\lim_{\lambda\to\infty}\int_0^1 f(x)\sin\lambda x\; dx=0.
\]

\item Si calcoli $\sum_{n=1}^{\infty}\frac 1{n^2}$ con una precisione
superiore ad un decimo.
\item Si risolva il seguente problema di Cauchy
\[
y'(x)=y(x)e^x\;;\quad y(0)=1.
\]
\item Si calcoli l'integrale
\[
\int_0^1\frac{x+3}{x^2+1}dx.
\]
\item Si trovi il polinomio di primo grado $p(x)$ tale che
\[
A(p):=\int_{-\pi}^{\pi}(\sin x -p(x))^2 dx
\]
\`e minimo.
\end{enumerate}
\end{document}
